% ======================================================================================
% Nom     : reports/latex/sections/10_conclusion.tex
% Rôle    : Section de conclusion du rapport.
% Auteur  : Maxime BRONNY
% Version : V1
% Cadre   : Réalisé dans le cadre du cours Techniques d'apprentissage artificiel
%           Master 1 Informatique Big Data
% Usage   : Fichier inclus dans main.tex lors de la compilation du rapport.
% ======================================================================================

\section{Conclusion}
\label{sec:conclusion}

\subsection{Bilan du travail réalisé}

Ce projet a permis de construire un pipeline d'apprentissage automatique
complet et reproductible pour la prédiction du risque de crédit bancaire :
chargement et nettoyage de 32\,581 dossiers, encodage des variables
catégorielles, découpage stratifié en trois ensembles ($70/15/15$), imputation
et normalisation sans fuite de données, entraînement de trois modèles de
familles différentes implémentés from scratch en Python, sélection des
hyperparamètres sur l'ensemble de validation, évaluation finale unique sur
l'ensemble de test, et génération automatique des métriques
(\texttt{results/metrics.json}) et des figures. L'ensemble s'exécute en une
seule commande et une vingtaine de secondes, et la justesse des trois
implémentations a été validée par comparaison avec scikit-learn (écarts d'AUC
inférieurs ou égaux à $5 \times 10^{-4}$).

\subsection{Réponse à la problématique}

Oui, le risque de défaut est prédictible avec une fiabilité utile à partir des
seules caractéristiques du dossier, et la comparaison désigne clairement
\textbf{l'arbre de décision} comme le meilleur compromis : F1 de $0{,}817$,
AUC de $0{,}909$ et une précision de $0{,}977$ qui rend ses alertes presque
toujours fondées, devant le $k$-NN (F1 $0{,}683$) et la régression logistique
(F1 $0{,}626$). La réponse doit cependant être nuancée par le plafond commun
des trois modèles : environ $30\,\%$ des défauts restent indétectés, limite qui
semble tenir à l'information disponible dans les données plus qu'aux
algorithmes eux-mêmes. Un déploiement réel exigerait donc soit des variables
supplémentaires, soit l'acceptation explicite de ce risque résiduel.

\subsection{Intérêt de la migration vers Python}

La réécriture du projet, initialement développé en C, s'est révélée bien plus
qu'une traduction. Le passage à Python et NumPy a divisé la taille du code par
plusieurs fois et l'a rendu directement confrontable à une référence
(scikit-learn), ce que la version C ne permettait qu'imparfaitement. Surtout,
l'effort économisé sur la gestion de la mémoire a pu être réinvesti dans le
protocole expérimental : ajout d'un ensemble de validation dédié (dont
l'apport est mesurable : environ douze points de F1 gagnés sur la régression
logistique par le seul réglage du seuil), ajout d'une troisième méthode, et
couverture par des tests automatisés. La compréhension fine des algorithmes
n'a rien perdu au change, puisque chaque méthode reste implémentée from
scratch, équations à l'appui.

\subsection{Bilan personnel d'apprentissage}

% Bilan rédigé à la première personne - à ajuster avec tes propres mots si tu
% le souhaites.
Ce projet m'a d'abord appris ce que « implémenter from scratch » exige
réellement : écrire une sigmoïde, c'est facile ; écrire une sigmoïde qui ne
déborde pas numériquement, une recherche de seuil de Gini qui tient en quelques
opérations vectorisées, ou un $k$-NN qui ne sature pas la mémoire, c'est là que
se loge la vraie compréhension des algorithmes. J'ai ensuite mesuré à quel
point le protocole expérimental compte autant que les modèles : la distinction
validation/test, que je percevais comme une convention, a produit sous mes yeux
un gain de douze points de F1 et détecté un début de surapprentissage. Enfin,
travailler sur des classes déséquilibrées m'a définitivement guéri de lire une
accuracy au premier degré : c'est la matrice de confusion, lue avec les coûts
métier en tête, qui dit ce qu'un classifieur vaut vraiment.
